\chapter{Chills}

\section*{Definition}
A feeling of shivering or trembling, often occurring with fever as the body changes its temperature set point. Chills can also occur without fever.

\section*{When to See a Doctor}
See your doctor if:
\begin{itemize}
  \item Chills persist, recur, or occur without an obvious cause, especially with fever or other illness
  \item Seek urgent care for confusion, severe weakness, difficulty breathing, a stiff neck, a purple rash, or signs of shock.
\end{itemize}

\section*{How It Develops}
During a fever, the brain may raise the body's temperature set point. Blood vessels narrow and muscles shiver to generate heat; chills can also reflect cold exposure, medicines, or other conditions.

\section*{Causes}
\begin{itemize}
  \item Bacterial infection
  \item Malaria
  \item Drug reaction
  \item Post-operative state
  \item Heat stroke recovery
\end{itemize}

\section*{Related Symptoms}
\begin{itemize}
  \item Fever
  \item Sweating
  \item Malaise
  \item Headache
  \item Muscle aches
\end{itemize}

\section*{Medical Evaluation}
\begin{itemize}
  \item History includes fever duration, pattern (intermittent, remittent, continuous), associated symptoms (chills, rigors, sweats), and exposure history.
  \item Vital signs assess fever severity and hemodynamic stability; complete physical examination seeks a source.
  \item Laboratory tests, cultures, and imaging are selected according to the examination, severity, likely source, and relevant exposures; they are not all needed for every episode.
  \item Imaging (chest X-ray, abdominal ultrasound, CT) is guided by clinical findings and suspected source.
\end{itemize}

\section*{Treatment Options}
\begin{itemize}
  \item Fever-reducing medications (acetaminophen, NSAIDs) provide symptomatic relief but do not treat the underlying cause.
  \item Antibiotics based on the most likely cause may be indicated for suspected bacterial infection after appropriate cultures are obtained.
  \item Antiviral, antifungal, or antimalarial therapy is directed at specific pathogens based on exposure and testing.
  \item Source control (drainage of abscess, removal of infected device) is essential for focal infections.
\end{itemize}

\section*{Self-Care and Home Management}
\begin{itemize}
  \item Keep warm and dry, and change damp clothing.
  \item Treat the cause of chills rather than repeatedly suppressing the symptom; seek care for chills with fever or feeling seriously unwell.
  \item Record temperature, timing, exposures, medicines, and associated symptoms.
  \item Seek urgent care for chills with severe illness, confusion, breathing difficulty, stiff neck, or persistent high fever.
\end{itemize}

\section*{References}
\begin{thebibliography}{99}
\bibitem{chills:ref1} Ogoina D. ``Fever, fever patterns and diseases called 'fever'--a review.'' J Infect Public Health. 2011;4(3):108-124.
\bibitem{chills:ref2} Plaisance KI, Mackowiak PA. ``Antipyretic therapy: physiologic rationale, diagnostic implications, and clinical consequences.'' Arch Intern Med. 2000;160(4):449-456.
\bibitem{chills:ref3} Dinarello CA, Porat R. ``Fever and the inflammatory response.'' In: Harrison's Principles of Internal Medicine. 21st ed. McGraw-Hill; 2022.
\bibitem{chills:ref4} Young PJ, Saxena M, Beasley R, et al. ``Fever and hypothermia in the critically ill.'' Crit Care. 2012;16(1):104.
\bibitem{chills:ref5} Niven DJ, Laupland KB, Tabah A, et al. ``Diagnosis and management of fever in the intensive care unit.'' Intensive Care Med. 2020;46(9):1733-1745.
\bibitem{chills:ref6} Mackowiak PA. ``Fever in adults.'' UpToDate. Wolters Kluwer; 2025.
\end{thebibliography}
